\documentclass[11pt,a4paper]{article}

\usepackage[margin=0.85in]{geometry}
\usepackage{setspace}
\usepackage{fontspec}
\usepackage{xeCJK}
\setCJKmainfont{Microsoft YaHei}
\usepackage{hyperref}
\usepackage{enumitem}
\usepackage{titlesec}

\hypersetup{colorlinks=true, urlcolor=blue, linkcolor=black}
\setlist[itemize]{leftmargin=*, itemsep=0.15em, topsep=0.2em, parsep=0pt}
\setlength{\parindent}{0pt}
\setlength{\parskip}{0.35em}
\pagestyle{empty}

\titleformat{\section}{\large\bfseries\uppercase}{\thesection}{0em}{}[\titlerule]
\titlespacing*{\section}{0pt}{1.1em}{0.55em}

\begin{document}

\begin{center}
{\LARGE\bfseries HSU, Ke-Cheng (Clio)}\\[0.4em]
{\large 許可證}\\[0.55em]
(852) 6840 9753 \quad\textbar\quad
\href{mailto:khsu@connect.ust.hk}{khsu@connect.ust.hk} \quad\textbar\quad
\href{https://cliohsu1997.github.io}{https://cliohsu1997.github.io} \quad\textbar\quad
\href{https://github.com/cliohsu1997}{github.com/cliohsu1997}
\end{center}

\vspace{0.4em}
\textbf{BUSINESS ADDRESS}\\
Department of Economics, Lee Shau Kee Business Building\\
The Hong Kong University of Science and Technology, Clear Water Bay, Kowloon, Hong Kong

\vspace{0.3em}
\textbf{DESIRED RESEARCH AND TEACHING FIELDS}\\
Industrial Organization, Applied Microeconomics, Empirical Industrial Organization

\section*{Education}
\textbf{Hong Kong University of Science and Technology}\\
Ph.D.\ in Economics \hfill 2020 -- 2027 (expected)\\
Research interests: Industrial Organization and Applied Microtheory.

\vspace{0.35em}
Master of Philosophy in Economics \hfill 2020 -- 2022

\vspace{0.35em}
Bachelor of Science in Mathematics and Economics \hfill 2015 -- 2019\\
Minor in Actuarial Mathematics

\section*{Working Papers}
\begin{itemize}
  \item ``Designing Return Policy in E-Commerce: The Role of Excess Return and Surplus Extraction'' \textit{(Job Market Paper)}\\
  \textit{Abstract:} Return policy design is central to social welfare, yet its welfare effects are ambiguous when policies jointly shape purchases, prices, and return behavior. Policy effects operate through both return discipline and market participation, and the two margins can pull in opposite directions. We estimate a structural model of purchase and return decisions using Taobao transaction data under RRSI and CRSI, and evaluate return-penalty discount (RPD) alternatives. Return policy should be tailored to category margins and return-handling costs rather than imposed uniformly.\\
  Presentations: HKUST Business School PhD Conference; CEP Brownbag, HKUST; 中国经济学年会夏季论坛$^{\ast}$ ($^{\ast}$ scheduled)

  \item ``Product Design with Heterogeneous Search Costs'' \textit{(MPhil thesis)}, with AU Pak Hung\\
  \textit{Abstract:} Building on Wolinsky's (1986) sequential search model with differentiated products, this paper introduces search cost heterogeneity to analyze how firms leverage observable product designs to influence consumer search behavior. In equilibrium, low-search-cost consumers tend to seek niche products, while higher-search-cost consumers gravitate toward standardized products. Prices may respond asymmetrically to reductions in search costs depending on where the shock occurs in the search-cost distribution.\\
  Presentations: APIOC 2023; AMES 2023, NTU Singapore
\end{itemize}

\section*{Work in Progress}
\begin{itemize}
  \item ``Too Much Unsold Stock?: Environmental Consequences of Price Commitment'', with Kohei Kawaguchi and Kei Kawai
\end{itemize}

\section*{Professional Experience}
\textbf{Research}\\
Research Assistant for Prof.\ Benjamin Bernard, National Taiwan University \hfill 2020

\vspace{0.35em}
\textbf{Teaching}\\
Lecturer, Math Camp for MSc in Economics, HKUST \hfill 2026\\
Teaching Assistant Coordinator, Business School, HKUST \hfill 2025 -- 2026\\
Teaching Assistant, Department of Economics, HKUST \hfill 2020 -- 2025\\
Courses: Macroeconomics; Industrial Organization \& Competitive Strategy; Economics of Uncertainty and Information; Microeconomic Theory I; Microeconomic Theory II

\section*{Fellowships and Awards}
Postgraduate Studentship from the HKUST, 2020 -- present\\
Red Bird PhD Award, 2022\\
Academic Achievement Medal, 2019\\
University Scholarship (730K HKD), 2015 -- 2019

\section*{References}
Available upon request.

\end{document}
